\begin{tikzpicture}[
  font=\scriptsize,
  panah/.style={-{Stealth[length=2mm]}, thick},
  modul/.style={draw, rounded corners, fill=blue!8, minimum width=1.9cm,
    minimum height=0.6cm, text height=1.6ex, text depth=0.4ex},
  luar/.style={draw, rounded corners, fill=gray!35, minimum width=1.9cm,
    minimum height=0.6cm, align=center},
  basis/.style={draw, cylinder, shape border rotate=90, aspect=0.22,
    fill=gray!15, minimum width=1.6cm, minimum height=0.7cm, font=\tiny,
    align=center},
  akses/.style={{Stealth[length=1.5mm]}-{Stealth[length=1.5mm]}, gray!60!black},
  tersembunyi/.style={-{Stealth[length=2mm]}, thick, dashed, red!70!black},
  ket/.style={font=\tiny, gray!60!black}
]
  % Bagian atas: kondisi sekarang, semua modul membaca dan menulis tabel yang sama.
  \node[anchor=west, font=\scriptsize\bfseries] at (-0.2,7.15)
    {Sekarang: semua bagian berbagi tabel yang sama};
  \node[modul] (p1) at (1.0,6.3) {Pendaftaran};
  \node[modul] (p2) at (3.5,6.3) {Rekam medis};
  \node[modul] (p3) at (6.0,6.3) {Farmasi};
  \node[modul] (p4) at (8.5,6.3) {Penagihan};
  \node[modul] (p5) at (11.0,6.3) {Klaim BPJS};
  \node[basis, minimum width=3.6cm, minimum height=1.0cm] (db) at (6.0,4.55)
    {\textbf{Basis data SIMRS}\\tabel pasien, kunjungan,\\resep, tagihan};
  \draw[akses] (p1.south) -- (db.north west);
  \draw[akses] (p2.south) -- (db.north);
  \draw[akses] (p3.south) -- (db.north);
  \draw[akses] (p4.south) -- (db.north);
  \draw[akses] (p5.south) -- (db.north east);
  \node[luar] (apotek) at (11.0,4.55) {Aplikasi apotek\\
    {\tiny\mdseries di luar SIMRS}};
  \draw[tersembunyi] (apotek.west) -- node[ket, above, red!70!black]
    {baca langsung} (db.east);
  \node[draw, rounded corners, fill=red!22, text width=3.0cm, align=center,
    font=\tiny] at (1.4,4.55)
    {Satu kolom di tabel pasien diubah: enam bagian ikut diuji dan dirilis
    bersama};

  % Garis pemisah.
  \draw[dashed, gray] (-0.2,3.45) -- (12.2,3.45);

  % Bagian bawah: kondisi target, tiap modul memiliki datanya sendiri.
  \node[anchor=west, font=\scriptsize\bfseries] at (-0.2,3.0)
    {Target: tiap modul memiliki datanya, modul lain bertanya lewat antarmuka};
  \node[modul, fill=green!10] (q1) at (1.0,2.15) {Pendaftaran};
  \node[modul, fill=green!10] (q2) at (3.5,2.15) {Rekam medis};
  \node[modul, fill=green!10] (q3) at (6.0,2.15) {Farmasi};
  \node[modul, fill=green!10] (q4) at (8.5,2.15) {Penagihan};
  \node[modul, fill=green!10] (q5) at (11.0,2.15) {Klaim BPJS};
  \node[basis] (d1) at (1.0,0.85) {data\\pendaftaran};
  \node[basis] (d2) at (3.5,0.85) {data\\klinis};
  \node[basis] (d3) at (6.0,0.85) {data\\resep};
  \node[basis] (d4) at (8.5,0.85) {data\\tagihan};
  \node[basis] (d5) at (11.0,0.85) {data\\klaim};
  \foreach \nomor in {1,...,5} {
    \draw[akses] (q\nomor.south) -- (d\nomor.north);
  }
  \draw[panah] (q2) -- (q1);
  \draw[panah] (q3) -- (q2);
  \draw[panah] (q4) -- (q3);
  \draw[panah] (q5) -- (q4);
  \node[ket] at (2.25,2.62) {antarmuka};
  \node[ket] at (4.75,2.62) {antarmuka};
  \node[ket] at (7.25,2.62) {antarmuka};
  \node[ket] at (9.75,2.62) {antarmuka};
  \node[draw, rounded corners, fill=green!10, text width=9.0cm, align=center,
    font=\tiny] at (6.0,-0.2)
    {Satu kolom data pendaftaran diubah: cukup modul pendaftaran yang diuji,
    selama antarmukanya tetap};
\end{tikzpicture}
